\section{Ecce Ancilla Domini}

All magic is the putting into practice of the maxim: the subtle rules the dense, writes Tomberg. He quotes the words of Mary to the angel Gabriel at the head of the chapter. The empress symbolizes the super-consciousness ruling the consciousness (which rules force, which rules matter). The crown on the head of the empress represents the super-consciousness ruling consciousness. Papus defines magic as ``the application of the strengthened human will to accelerate the evolution of the living forces of Nature'', while Peladan defines magic as ``the art of the sublimation of man''. Magic is also ``the science of Love''. The attentive reader should think of Dante, because the stars are moved by Love, and the stars rule us, to begin with, as intermediaries. The aim is accelerated evolution, the means are Love, symbolized by the Eagle and the sceptre.

Valerie Flint has done good research\footnote{\url{http://www.amazon.com/Rise-Magic-Early-Medieval-Europe/dp/0691001103/ref=sr\_1\_1?ie=UTF8&qid=1410577722&sr=8-1&keywords=valerie+flint+magic}} on Christianity's unique relation to magic during the Middle Ages: Christianity baptized and used a great portion of it, for several reasons, in a traditional manner. She argues that Christianity succeeded in large measure because her priests were more effective at dealing with curses, disease, and moral darkness than were the pagan peers, and also because they learned what was best from old Europe. I encourage people to consider if this may need to become so once again.

Redemption, he goes on to say, requires the synergy of two wills: the divine and the human. Therefore, when Peter heals a man (Acts 9:32-34), the involvement of Peter is not a superfluity or accident. The healing required Peter; nevertheless, it is true that it came completely from God as well. Tomberg refers this mystery to the Incarnation, and it requires deep meditation. Sacred magic renders service, not merely to what is above, but to that which is below. Thus, it is not merely the vulgar connotation of ``magic'' that is of interest, but thaumaturgy. It is first necessary to contact the Divine (mysticism), then it is necessary to decode this in the understanding (gnosis), prior to the putting into action of that which is ``given''. The magus becomes a link between that which is higher, and aims to heal, and that which is lower, and requires assistance. The mage is crucified, but the magus is not a martyr: rather, he is a thaumaturgist\footnote{\url{http://en.wikipedia.org/wiki/Iamblichus}}. The aim of magic is ``liberated action''. Tomberg notes that had he not visited the Chapel of the Holy Blood at Bruges\footnote{\url{http://en.wikipedia.org/wiki/Basilica\_of\_the\_Holy\_Blood}}, he could not have written this chapter.

What does this mean? Cologero has written over and over that books do not a conversion make, let alone an internet career. Religion itself, to be vivified and vivifying, has to be founded in mysticism (faith), illumined by gnosis (hope) and actuated by sacred magic (love). It grows cold without faith, clouds without gnosis, and is impotent without magic.

And yet, the books are somehow necessary\footnote{\url{http://en.wikipedia.org/wiki/The\_Philobiblon}}. Who among us can say that, without books, we would pursue divine study with more acumen? We receive this learning, however, at second and third hand, all the same, and all the more. This is one of the meanings of Tradition. A man speaks, a disciple writes, a reader learns, a teacher passes on, a learner assimilates, a wise man remembers, and the flame passes down. It is a fact that super-consciousness guides the story of the Cosmos, and that in the Ashkaric records (which Tomberg takes to exist as a matter of fact), there is a living proof that even the Eternal itself keeps records, records events, remembers the passing human simulation, a simulation with a divine end and point.

There is a tendency in humanity not to take itself seriously enough, or to take itself seriously in a flat and horizontal way (and this split is a typical tendency towards dichotomy and flattened, binary thought). The ``cure'' for this proceeds along the lines of divine magic. Just as humans have a variety of means and methods for getting over physical illness, so do humans find a variety of means for getting over moral illness. This ``divine attractor'' (which can, quite frankly, be judged by its fruits) draws a great many parallel lines through a great many planes and worlds. These are curative: that is to say, they are (in Mouravieff's terms) ``B'' influences. It is up to the synergy of the human will with the divine to both discern and cooperate with them. \emph{That is all that we can do}. Initially, without Him, we can do nothing. It is only farther down the road that there ceases to be a division between the divine and human will in terms of friction and resistance.

If you reach for something, and fail, and yet that something stays beside you and gives itself to you, you have experienced the divine magic, because the point was to see that you could not possess it, it had to be given to you. And who else gives it to you but the Friend above''. In the Psalms, it teaches expressly that ``he gives gifts to those who are below'' (in Hades). How much more to those who pursue the disciplines of ``accelerated evolution''? This is a rigorously mathematical certainty, from the lesser to the greater. Glory to God for His long suffering (an ancient invocation for healing in any situation), the argument from the greater to the lesser is also valid! If He loves the righteous, will He not also seek out the last sheep? Who can possibly understand this? Why should God be operative in both directions, from below towards the ``above'', and from above towards the ``below''?

This is because God is Love: He is the energy that moves the universe, the stars, which moves us. But of course, we are meant to move the stars. The one way to do this, is to achieve the energy of being in the state of Love, mistress and empress of magic.


\flrightit{Posted on 2014-09-13 by Logres}